
\documentclass[12pt,a4paper]{article}

% 中文支持
\usepackage[UTF8]{ctex}

% 数学公式
\usepackage{amsmath,amssymb}

% 图片
\usepackage{graphicx}

% 页面设置
\usepackage{geometry}
\geometry{margin=2.5cm}

\title{Digital Photogrammetry Notes}
\author{PotatoPi}
\date{2026 Fall}

\begin{document}

\maketitle

\section{引入}
\subsection{关于数字摄影测量}
\subsection{影像、采样}

\section{重采样}
\subsection{双线性插值}
\subsection{双三次卷积}
\subsection{最邻近像元法}

\section{特征提取}
\subsection{点特征提取}
\subsubsection{Moravec算子}
\subsubsection{Forstner算子}
\subsubsection{Harris算子}

\subsection{线特征提取}
\subsubsection{一阶差分算子}
\subsubsection{二阶差分算子}

\subsection{Hough变换}

\section{摄影测量解析基础}
\subsection{内定向}
\subsection{核线几何关系解析}

\section{影像匹配}
\subsection{基本概念}
\subsection{匹配算法}

\section{最小二乘影像匹配}

行内公式：$E=mc^2$

独立公式：

\[
    \frac{\partial f}{\partial x}
    =
    \lim_{h\to 0}
    \frac{f(x+h)-f(x)}{h}
\]



\[
A =
\begin{bmatrix}
1 & 2 \\
3 & 4
\end{bmatrix}
\]

\end{document}
